\documentclass[11pt]{article}

% ---------- Packages ----------
\usepackage[a4paper,margin=1in]{geometry}
\usepackage{amsmath,amssymb,amsthm,mathtools}
\usepackage{hyperref}
\usepackage{enumitem}
\usepackage{booktabs}
\usepackage{array}

% ---------- Theorem environments ----------
\newtheorem{theorem}{Theorem}[section]
\newtheorem{proposition}[theorem]{Proposition}
\newtheorem{lemma}[theorem]{Lemma}
\newtheorem{corollary}[theorem]{Corollary}

\theoremstyle{definition}
\newtheorem{definition}[theorem]{Definition}
\newtheorem{example}[theorem]{Example}
\newtheorem{remark}[theorem]{Remark}
\newtheorem{question}[theorem]{Question}

% ---------- Commands ----------
\newcommand{\Ccal}{\mathcal{C}}
\newcommand{\St}{\mathrm{St}}
\newcommand{\Clop}{\mathrm{Clop}}
\newcommand{\pure}{\mathrm{pure}}
\newcommand{\CO}{\mathrm{CO}}
\newcommand{\tr}{\mathrm{tr}}

% ---------- Title ----------
\title{The Extension Problem for Orthomodular\\
Observation Algebras}
\author{Patrick S.\ Mahon}
\date{Research note --- June 2026}

\begin{document}
\maketitle

% ==================================================================
\section{Background}
\label{sec:background}
% ==================================================================

All measurement is finite.  An experimenter performs finitely many
trials, records finitely many digits, distinguishes finitely many
outcomes.  Yet the mathematical frameworks we use to organise
empirical knowledge --- probability theory, quantum mechanics,
statistical mechanics --- rest on infinite structures:
$\sigma$-algebras, Hilbert spaces, continuous symmetry groups.
The passage from finite observation to infinite structure is not
automatic; it requires a commitment that no finite body of
evidence can compel.

In probability theory, the locus of this commitment is
\emph{countable additivity}.  Given a Boolean algebra
$\mathcal{E}$ of events and a finitely additive charge
$\ell : \mathcal{E} \to [0,1]$ with $\ell(\Omega) = 1$,
the condition
\begin{equation}\label{eq:sigma-additivity}
  E_n \downarrow \varnothing
  \quad\Longrightarrow\quad
  \ell(E_n) \to 0
\end{equation}
--- that mass cannot persist along a vanishing sequence of events
--- is not derivable from any finite collection of observations.
De Finetti~\cite{deFinetti1974} argued that only finite additivity
is empirically grounded; Kolmogorov~\cite{Kolmogorov} adopted
$\sigma$-additivity as an axiom and built modern probability on
it.  The question of \emph{when} and \emph{why} the passage from
one to the other is justified has never been fully resolved.

In quantum theory, the same tension appears in a different guise.
The algebra of propositions about a quantum system is not Boolean:
incompatible observables cannot be simultaneously determined.  The
closed subspaces of a Hilbert space $H$ form an orthomodular
lattice $L(H)$ --- a structure in which complementation is
well-defined but distribution of meet over join may fail:
\begin{equation}\label{eq:non-distributive}
  a \wedge (b \vee c) \;\neq\; (a \wedge b) \vee (a \wedge c)
  \qquad\text{in general.}
\end{equation}
A \emph{state} on such a lattice is a function
$s : L(H) \to [0,1]$ satisfying
\begin{equation}\label{eq:state-def}
  s(\mathbf{1}) = 1,
  \qquad
  s(a \vee b) = s(a) + s(b) \quad\text{whenever } a \perp b.
\end{equation}
The non-distributivity means that $s$ is additive only on
orthogonal pairs, not on arbitrary finite unions.  The classical
machinery for extending finitely additive assignments to
$\sigma$-additive ones --- Carath\'{e}odory's outer-measure
construction, which requires a Boolean algebra of sets --- does
not apply.

For the specific lattice $L(H)$, this difficulty is resolved by a
remarkable rigidity.  Gleason's theorem~\cite{Gleason1957} shows
that, when $\dim H \geq 3$, every state has the form
\begin{equation}\label{eq:gleason}
  s(P) \;=\; \tr(\rho\, P)
\end{equation}
for a unique density operator $\rho$, and is therefore
automatically $\sigma$-additive.  The geometry of Hilbert space is
so constraining that the problem of the infinite simply does not
arise --- states have no room to misbehave.

But Hilbert space is very special.  Orthomodular lattices arise in
many other settings --- the projection lattices of operator
algebras, the logics of quantum systems with superselection rules,
and the abstract lattices studied in quantum logic --- and for
these, no analogue of Gleason's rigidity is known.  The question
that probability theory and quantum theory share --- when does
finite coherence force countable behaviour? --- remains open in
this broader setting, and is the subject of the present note.

\medskip

The programme ``Structure from Observation''~\cite{mahon2025}
approaches this question from a particular direction: rather than
beginning with a pre-given space $\Omega$ of outcomes and asking
which $\sigma$-algebras it supports, it begins with the
\emph{observations themselves} --- a directed system of event
algebras $\mathcal{E}_i$ with compatible charges $\ell_i$ --- and
asks what the directed structure alone determines.

Paper~I answers this for Boolean event algebras.  A compatible
family of finitely additive charges extends to a $\sigma$-additive
probability measure if and only if collective exhaustion (CE) holds
--- condition~\eqref{eq:sigma-additivity} applied uniformly across
the directed system.  Two constructions realise the extension:
\begin{itemize}[noitemsep]
  \item \textbf{Carath\'{e}odory:} assumes CE; produces a measure
    on $(\Omega,\, \sigma(\Ccal))$ directly.
  \item \textbf{Stone:} assumes nothing; produces a Baire measure
    on the Stone space $\St(\Ccal)$ unconditionally.
\end{itemize}
Under discriminability, the two agree.  Three questions orient the
broader programme:
\begin{enumerate}[label=(\arabic*)]
  \item What must be added to observation to get structure?
  \item Reconstruction without points: what structure is already
    present in the relational description --- contexts, directed
    refinement --- before point-spaces are assumed?
  \item Local-to-global extension: what conditions on the
    \emph{site} control whether local data extend globally?
\end{enumerate}
Paper~I resolves~(2) and~(3) for Boolean algebras.  The present
note asks whether any of this survives when the event algebra is
orthomodular but not distributive.

Epperson and Zafiris~\cite{EppersonZafiris2013} advocate a
\emph{relational realist} ontology in which structure is
constituted by relations among contexts, not by points in a
pre-given space.  An orthomodular lattice $A$ is the natural
algebraic expression of this: its elements are propositions, its
partial order is entailment, and its non-distributivity reflects
the impossibility of assigning simultaneous truth values to all
propositions.  The dual space $S_0(A)$ of
McDonald--Bimb\'{o}~\cite{McDonaldBimbo2023} is built from
\emph{filters} --- consistent collections of propositions --- not
from points of a pre-existing space.  This is question~(2) applied
to non-Boolean algebras: what structure does the relational
description already contain?

In the Boolean case, one might hope that the passage from finite to
countable additivity could be understood as a \emph{sheaf condition}
--- local states gluing to a global $\sigma$-additive measure.
Biesel~\cite{Biesel2024} shows this hope is disappointed:
\begin{itemize}[noitemsep]
  \item Finitely additive charges form a sheaf for finite covers.
  \item $\sigma$-additive measures form a sheaf for countable
    covers.
  \item The step between is the \emph{definition} of
    $\sigma$-additivity, not a structural theorem.
\end{itemize}
This is the boundary at which the present problem begins.

Paper~II~\cite{MahonPaperII2026} introduces three positions on the
relationship between observation and reality, each a constraint on
where the dual-space measure lives:
\begin{description}[style=nextline, font=\normalfont\bfseries,
    labelwidth=3.5cm, leftmargin=4cm]
  \item[Empirical adequacy (EA)]
    The dual-space measure agrees with all observations.
  \item[Probabilistic realism (PR)]
    The measure descends to a realisation space.
  \item[Value-definite realism (VDR)]
    A global two-valued homomorphism exists.
\end{description}
These are nested: $\text{VDR} \Rightarrow \text{PR}
\Rightarrow \text{EA}$.  For Boolean algebras, all three are
commensurable --- every EA-theory admits a VDR-completion.  For
OMLs with $\dim \geq 3$, Gleason's
theorem~\eqref{eq:gleason} rescues PR, but
Kochen--Specker~\cite{KochenSpecker1967} blocks VDR.  The
extension problem sits at the $\text{EA} \to \text{PR}$ transition:
can we pass from a state on the algebra (EA) to a
$\sigma$-additive measure on the dual space (PR)?  If so, under
what conditions --- and is descent then a free choice (as in the
Boolean case) or a structural consequence of the algebra and state
together?

Van Fraassen's distinction~\cite{vanFraassen1980} between
empirical adequacy and structural commitment, which Paper~I makes
precise via the Stone construction, thus acquires a sharper form
in the OML setting: the Stone route is no longer available, the
Carath\'{e}odory route is blocked by non-distributivity, and the
passage from coherent local data to global probability ---
question~(3) --- requires genuinely new ideas.

% ==================================================================
\section{The problem}
\label{sec:problem}
% ==================================================================

Let $A$ be an orthomodular lattice (OML) and let $s : A \to [0,1]$
be a state~\eqref{eq:state-def}.  The McDonald--Bimb\'{o}
duality~\cite{McDonaldBimbo2023} associates to $A$ a compact
topological space
\[
  S_0(A) \;=\; \bigl(F(A),\, \subseteq,\, \perp_A,\, P(A),\,
  T(S)\bigr),
\]
where $F(A)$ is the set of filters on $A$ and
$P(A) \subseteq F(A)$ is the set of prime filters.  For each
$a \in A$, set
\[
  h(a) \;=\; \{x \in F(A) : a \in x\}.
\]
The map $h$ is an OML isomorphism from $A$ onto the
$\perp$-stable clopen sets $\CO(S_0(A))^\dagger$.  Define a set
function $\mu$ on these clopens by
\begin{equation}\label{eq:mu-def}
  \mu\bigl(h(a)\bigr) \;=\; s(a).
\end{equation}

\begin{question}[Extension]\label{q:extension}
Under what conditions on the OML $A$ does $\mu$ extend to a
$\sigma$-additive measure on the Baire (or Borel) $\sigma$-algebra
of $S_0(A)$?
\end{question}

This is the direct OML analogue of the Stone extension theorem
established for Boolean algebras in~\cite{mahon2025}.

In the Boolean case, the extension is automatic.  The clopens
$\Clop(\St(\Ccal))$ form a Boolean algebra, the state $s$
restricts to a finitely additive charge on this algebra, and the
Carath\'{e}odory extension theorem (or Choksi's projective-limit
argument~\cite{choksi1958}) produces the $\sigma$-additive
extension.  Compactness of the Stone space gives
$\sigma$-subadditivity for free.

In the OML case each of these steps breaks:
\begin{enumerate}[label=(\roman*)]
  \item $\CO(S_0(A))^\dagger$ is an OML, \emph{not} a Boolean
    algebra.
  \item The state $s$ is additive on orthogonal pairs only, not on
    arbitrary finite unions.
  \item \textbf{Carath\'{e}odory does not apply} --- there is no
    Boolean algebra of sets on which to run the outer-measure
    construction.
  \item Compactness holds (McDonald--Bimb\'{o},
    Lemma~3.5~\cite{McDonaldBimbo2023}), but the
    $\sigma$-subadditivity argument requires orthogonality
    structure that the topology alone does not provide.
\end{enumerate}
The gap is sharp: ``finitely additive on a compact Boolean algebra
of clopens'' automatically extends; ``orthogonally additive on a
compact OML of clopen $\perp$-stable sets'' does \emph{not}.

For the prototypical case $A = L(H)$ with $\dim H \geq 3$,
Gleason's theorem~\eqref{eq:gleason} resolves the problem
completely: every state is induced by a density operator, and the
resulting measure is $\sigma$-additive on $L(H)$.  The extension
to $S_0(L(H))$ follows.  For general OMLs, no Gleason-type
representation is available.  Question~\ref{q:extension} asks for
the structural conditions that substitute for the Hilbert-space
geometry.

\begin{table}[ht]
\centering
\renewcommand{\arraystretch}{1.3}
\begin{tabular}{@{}l p{4.8cm} p{5.2cm}@{}}
\toprule
\textbf{Feature} & \textbf{Boolean (Paper~I)} & \textbf{OML (this problem)} \\
\midrule
Dual space
  & $\St(A)$ (ultrafilters)
  & $S_0(A)$ (all filters) \\
Distinguished subset
  & $\pure(\Omega)$, not canonical
  & $P(A)$, canonical \\
Clopens
  & Boolean algebra
  & OML (non-distributive) \\
Additivity of $s$
  & full (finite)
  & orthogonal pairs only \\
Extension
  & automatic (Carath\'{e}odory)
  & requires Gleason-type result \\
Realisation constrained?
  & No
  & Yes (Kochen--Specker~\cite{KochenSpecker1967}) \\
\bottomrule
\end{tabular}
\caption{Structural comparison of the extension problem.}
\label{tab:comparison}
\end{table}

\paragraph{Two axes: extension and descent.}
The passage from a state on $A$ to a measure on $S_0(A)$ splits into
two axes that the Boolean case fuses but the OML case separates.
Since $S_0(A)$ is a Stone space, its \emph{full} clopen algebra
$\CO(S_0(A))$ is Boolean, whereas the state lives only on the
$\perp$-stable clopens $\CO(S_0(A))^\dagger$ --- the OML, the image
of $h$.  These share only meet; join, complement, and bottom differ.
\begin{description}[style=nextline, font=\normalfont\bfseries,
    labelwidth=2.5cm, leftmargin=3cm]
  \item[Extension]
    Extend $\mu$ from the $\perp$-stable clopens to a finitely
    additive charge on the \emph{full} Boolean algebra of clopens.
    This requires consistent values on non-$\perp$-stable clopens
    --- unions $h(a) \cup h(b)$ for non-orthogonal $a, b$, which lie
    outside $\mathrm{im}(h)$ whenever $A$ is non-Boolean.  This is
    the Pt\'{a}k--Pulmannov\'{a} frontier; $\sigma$-additivity of
    the state does not help here.
  \item[Compactness]
    Once extension succeeds, the charge extends to a
    $\sigma$-additive Borel measure automatically.  Unconditional.
\end{description}
In the Boolean case (Paper~I), the extension axis is \emph{free}:
finite additivity on the clopen Boolean algebra plus compactness
gives the Stone measure unconditionally --- ``the first arrow is
free.''  In the OML case, \textbf{the extension axis is no longer
free} --- it is the open problem.  The non-distributivity that blocks
Carath\'{e}odory is exactly the obstruction to extending past the
orthogonal pairs.

\paragraph{Descent.}
$\sigma$-additivity of the state governs descent, not extension.
Even once $\mu$ extends to a measure $\hat\mu$ on $S_0(A)$, the
\emph{descent} question --- whether $\hat\mu$ concentrates on the
``physical'' points --- changes character:
\begin{description}[style=nextline, font=\normalfont\bfseries,
    labelwidth=2.5cm, leftmargin=3cm]
  \item[Boolean]
    Does $\hat\mu$ concentrate on $\pure(\Omega)$?  If and only if
    the charges are $\sigma$-additive.  The choice of realisation
    space $\Omega$ is unconstrained.
  \item[OML]
    Does $\hat\mu$ concentrate on $P(A)$?  Here $P(A)$ is part of
    the structure, not a free choice.  The Kochen--Specker
    obstruction~\cite{KochenSpecker1967} prevents simultaneous
    realisation of all filters; the Bub--Clifton
    theorem~\cite{BubClifton1996} shows that the state determines
    a maximal Boolean subalgebra of definite values.
\end{description}
In the OML setting, the algebra and state \emph{together}
constrain which filters are ``realised.''  Descent is not a
geometric commitment but a structural consequence.

\begin{remark}[Decomposition lives on the descent axis]
De Simone--Navara~\cite{DeSimoneNavara2001} decompose OMP states as
$s = s_\sigma + s_{\mathrm{wpfa}}$ (a $\sigma$-additive part and a
weakly purely finitely additive part), the orthomodular analogue of
Yosida--Hewitt.  By the axis split above, the wpfa component governs
the \emph{descent} axis --- it is the analogue of the purely finitely
additive part $\ell_p$ in the Boolean case --- and \emph{not} the
extension axis.  A natural-looking conjecture, ``$s$ extends iff
$s_{\mathrm{wpfa}} = 0$,'' is therefore mis-targeted: vanishing of
$s_{\mathrm{wpfa}}$ concerns whether the (already $\sigma$-additive)
state's measure concentrates on $P(A)$, while extension is blocked by
non-distributivity regardless of the wpfa component.  The
decomposition theory and the extension problem live on different
axes.

A caveat compounds this.  The McDonald--Bimb\'{o}
duality~\cite{McDonaldBimbo2023} is \emph{finitary}: the identity
$h(\bigvee_n a_n) = (\bigcup_n h(a_n))^{\perp\perp}$ holds for finite
joins but fails in general for countable joins, since finitary OML
homomorphisms need not preserve infinite suprema.  Making even the
descent argument rigorous --- the analogue of Paper~I's
``$\sigma$-additivity $\Leftrightarrow$ concentration'' --- requires
$\sigma$-completeness and continuity hypotheses the 2023 duality does
not supply, and is itself open.
\end{remark}

% ==================================================================
\section{What is known}
\label{sec:known}
% ==================================================================

In the Boolean case, the Kelley--Vladimirov--Pt\'{a}k
characterisation (Fremlin,
Theorem~391D~\cite{Fremlin3}) gives a complete algebraic answer:

\begin{theorem}[KVP]\label{thm:kvp}
A Boolean algebra $B$ carries a strictly positive
$\sigma$-additive measure if and only if $B$ is
\begin{enumerate}[label=(\alph*), noitemsep]
  \item Dedekind $\sigma$-complete,
  \item weakly $(\sigma,\infty)$-distributive, and
  \item chargeable (carries a strictly positive finitely additive
    measure).
\end{enumerate}
\end{theorem}

Each condition is algebraic and non-trivially constraining.  For
OMLs, the analogues fail or degenerate:
\begin{enumerate}[label=(\alph*), noitemsep]
  \item Dedekind $\sigma$-completeness makes sense but is very
    strong.
  \item Weak $(\sigma,\infty)$-distributivity has no clean OML
    analogue --- the definition relies on distributivity of
    $\bigwedge$ over $\bigvee$, which may fail orthomodularly.
  \item Chargeability has no known structural characterisation on
    OMLs; the existence of a strictly positive finitely additive
    measure cannot be read off from the lattice structure.
\end{enumerate}

The positive results that do exist for OMLs all require algebraic
structure rich enough to approximate Hilbert-space geometry:
\begin{description}[style=nextline, font=\normalfont\bfseries,
    labelwidth=4cm, leftmargin=4.5cm]
  \item[Gleason~\cite{Gleason1957}]
    $A = L(H)$, $\dim H \geq 3$.  Every state has the
    form~\eqref{eq:gleason}; $\sigma$-additivity is automatic.
  \item[Bunce--Wright~\cite{BunceWright1992}]
    $A$ the projection lattice of a JBW-algebra.
    $\sigma$-additivity via operator-algebraic structure.
  \item[Chetcuti--Dvure\v{c}enskij~\cite{ChetcutiDvurecenskij2005}]
    Lattice effect algebras with completeness conditions close to
    von~Neumann algebras.
\end{description}

The obstructive results show that the Boolean extension machinery
cannot be transplanted:
\begin{description}[style=nextline, font=\normalfont\bfseries,
    labelwidth=4cm, leftmargin=4.5cm]
  \item[Pt\'{a}k--Pulmannov\'{a}~\cite{PtakPulmannova1991}]
    If every unital subadditive measure on an orthomodular poset is
    a state, the poset must be Boolean.  This is the structural
    reason no KVP analogue exists for OMLs: conditions strong enough
    to force $\sigma$-additivity collapse the lattice back to a
    Boolean algebra.
  \item[Navara~\cite{Navara1992}]
    Regularity does \emph{not} force $\sigma$-additivity on general
    orthomodular posets.  The B\'{e}aver--Cook result is special to
    von~Neumann algebras.
  \item[Pt\'{a}k~\cite{Ptak1987}]
    Exotic OML state spaces exist; the obstruction is combinatorial
    (Greechie pasting), not from the centre.
\end{description}

Recent incremental work confirms the field is active but no one
has found the right condition:
Vor\'{a}\v{c}ek--Pt\'{a}k~\cite{VoracekPtak2023} study signed
measures on orthomodular posets;
Bure\v{s}ov\'{a}--Pt\'{a}k~\cite{BuresovaPtak2023} investigate
variations on regularity conditions; De~Simone--Navara study
Yosida--Hewitt-type decompositions for orthomodular posets.

The gap between ``too weak'' (allows non-$\sigma$-additive states)
and ``too strong'' (collapses to Boolean) appears structurally
robust.  This is not a problem where further reading will produce
the answer --- it requires a new idea.

All known exotic OML state spaces are built by \textbf{Greechie
pasting}~\cite{Wilce2009, Shultz1974, Ptak1987}.  Progress on the
extension problem likely requires either:
\begin{enumerate}[label=(\alph*)]
  \item a new pasting technique that controls $\sigma$-additivity,
    or
  \item an approach that bypasses the state space entirely (e.g.,
    categorical or
    topos-theoretic~\cite{DoeringIsham2008}).
\end{enumerate}

% ==================================================================
\section{Assessment}
\label{sec:assessment}
% ==================================================================

Question~\ref{q:extension} is a precise, well-formulated problem.
A positive answer would unify three results under a single
framework --- states on OMLs $\to$ measures on dual spaces $\to$
constrained descent:
\begin{enumerate}[label=(\roman*), noitemsep]
  \item Paper~I~\cite{mahon2025}: the Boolean case.
  \item Gleason's theorem~\cite{Gleason1957}: the $L(H)$ case.
  \item The Bub--Clifton theorem~\cite{BubClifton1996}: descent
    and realisation.
\end{enumerate}

\medskip\noindent
\textbf{Difficulty.}\enspace High.  The Pt\'{a}k--Pulmannov\'{a}
obstruction shows this is not merely ``find the right condition''
--- conditions strong enough to force extension destroy the
non-Boolean structure.

\medskip\noindent
\textbf{Novelty.}\enspace The formulation via McDonald--Bimb\'{o}
duality~\cite{McDonaldBimbo2023} appears to be new.  The
D\"{o}ring--Isham topos approach~\cite{DoeringIsham2008} addresses
related questions but via presheaves, not the filter-space duality.

\medskip\noindent
\textbf{Status.}\enspace Open but structurally resistant.  Not a
current active lead --- needs collaboration or a new idea.

\bibliographystyle{plain}
\bibliography{references_oml}
\end{document}
